\section{Redes Neuronales Artificiales}
Una red neuronal artificial es un modelo de computación basado en las redes neuronales 
biológicas. Una red neuronal está compuesta por nodos computacionales,
conocidos como neuronas, que forman conexiones y de forma colectiva aprenden para
optimizar su salida.

\subsection{La Neurona Artificial}
Las neuronas artificiales, inspiradas en las
neuronas biológicas, son nodos computacionales que tienen un valor de activación
similar a sus contrapartes biológicas. Esta activación es obtenida al recibir señales 
de entrada, que son combinan mediante una función
y producen un valor de salida. Las dos funciones principales de una neurona artificial
son la suma ponderada y la función de activación.

\begin{description}
    \item[Suma ponderada] Recibe los valores de salida de otras neuronas y los combina 
    considerando el peso de sus conexiones para obtener la entrada neta $z$ de la neurona.
    Está definida como:
    \begin{equation}
        \text{z} = \sum_{i=1}^{n}x_{i} w_{i} + b
    \end{equation}
    donde:
    \begin{description}
        \item[$x_{i}$] es el valor de salida de la $i$-ésima neurona de entrada.
        \item[$w_{i}$] es el peso asociado a la conexión de la $i$-ésima neurona de entrada.
        \item[$b$] es el valor del sesgo (bias).
        \item[$n$] es el número total de neuronas de entrada.
    \end{description}
    \item[Función de activación] Transforma la entrada neta $z$ en un valor de activación
    $a$ mediante una función no lineal:
    \begin{equation}
        a = f(z)
    \end{equation}
\end{description}
%% Mencionar bias

Estas dos funciones permiten modelar el comportamiento básico de una neurona artificial 
dentro de una red neuronal.
\begin{figure}[ht]
    \centering
    \includegraphics[width=0.8\textwidth]{diagrams/chapter2/neuron.pdf} 
    \caption{Modelo de la neurona artificial.}
    \label{fig:neurona}
\end{figure}

%% Agregar diagrama de una neurona

\subsection{Funciones de Activación}
Como se mencionó previamente, la neurona artificial transforma la entrada neta mediante
una función de activación. Estas funciones son fundamentales en las redes neuronales,
ya que introducen no linealidad en el modelo, permitiendo la representación de relaciones
complejas entre las variables de entrada y salida. En ausencia de funciones de activación
no lineales, la red se comportaría como una transformación lineal, independientemente del
número de capas.

%Entre las funciones de activación más utilizadas se encuentran las funciones sigmoides y 
%las funciones lineales por partes.

\subsubsection{Funcion Escalonada}
También conocida como \textit{función escalonada binaria}, es una función en la que la neurona 
es activada si la entrada de la neurona supera un umbral. La función escalonada está dada por:
\begin{equation}
    f(z) = \begin{cases}
        1 & \text{si } z \ge \theta \\
        0 & \text{si } z < \theta
    \end{cases}
\end{equation}
donde $\theta$ es el valor del umbral, comúnmente igual a $0$.
\medskip

La función escalonada es una de las funciones de activación más simples y puede ser 
utilizada en problemas de clasificación binaria. Sin embargo, presenta limitaciones 
importantes, ya que no es diferenciable y su gradiente es nulo en casi todo su dominio. 
Esto dificulta el uso de algoritmos de retropropagación durante el proceso de aprendizaje, lo que limita su aplicación en redes neuronales modernas.

\subsubsection{Función Lineal}
En el contexto de redes neuronales, la función lineal utilizada como función de activación 
corresponde típicamente a la función identidad:
\begin{equation}
    f(z) = z
\end{equation}
A diferencia de la función escalonada, esta función es diferenciable y su gradiente es 
constante, por lo que los pesos pueden ser actualizados mediante algoritmos de 
retropropagación.

Sin embargo, presenta una limitación importante: si todas las funciones de activación de 
una red son lineales, entonces la red es equivalente a una única transformación lineal, 
lo que impide modelar relaciones no lineales y reconocer patrones complejos.

Por esta razón, la función lineal se utiliza típicamente en la capa de salida de modelos 
de regresión.

\subsubsection{Función Sigmoide}
La función sigmoide logística es una función no lineal que transforma la entrada en un valor
dentro del intervalo $(0,1)$, definida como:
\begin{equation}
    f(z)=\frac{1}{1+e^{-z}}
\end{equation}
Es una de las funciones de activación más utilizadas ya que mapea valores al rango $(0,1)$
al prevenir valores de activación muy altos. Presenta una limitación importante conocida
como \guillemotleft gradiente desvanescente\guillemotright, en la que valores hacia ambos
extremos de la función producen gradientes muy pequeños, lo que produce cambios poco significativos
que dificultan la retropropagación.
\medskip

Adicionalmente, al no estar centrada en cero, las activaciones de la función sigmoide 
son siempre positivas, lo que puede generar actualizaciones correlacionadas (los gradientes
tienen el mismo signo) en los pesos y ralentizar el proceso de optimización.
\medskip

La función sigmoide se utiliza comúnmente en la capa de salida de redes neuronales
para problemas de clasificación binaria, ya que transforma la salida en un valor
en el intervalo $(0,1)$, el cual puede interpretarse como una probabilidad de
pertenencia a una clase.

\subsubsection{Función Tangente Hiperbólica}
La función hiperbólica tangente, también conocida solo como función $\tanh$, es una
función de activación no lineal similar a la función sigmoide. 
Sin embargo, a diferencia de esta, la función tangente es simétrica respecto al origen.
Está definida como:
\begin{equation}
    f(z) = \tanh{z} = \frac{e^{z}-e^{-z}}{e^{z}+e^{-z}}
\end{equation}
La función tangente hiperbólica produce valores en el rango $(-1,1)$, por lo que suele ser
preferida frente a la función sigmoide, ya que sus salidas están centradas en
cero, lo que permite obtener derivadas derivadas de mayor magnitud en la región cercana 
a cero.
Sin embargo, aún presenta el problema del desvanecimiento del gradiente debido a la
saturación de la función para valores grandes de la entrada.

\subsubsection{Función ReLU}
La función ReLU, por sus siglas en inglés \textit{Rectified Linear Unit}, es una función 
de activación lineal por partes que coincide con la función identidad para valores positivos y toma el 
valor de cero en cualquier otro caso. También conocida como la función rampa, se define como:
\begin{equation}
    f(z) = \max(0,z)
\end{equation}
Una de sus principales ventajas es su simplicidad computacional, ya que su gradiente toma valores constantes 
$0$ o $1$, lo que facilita el cálculo durante la retropropagación
Esto ha contribuido a su amplio uso en redes neuronales profundas.
\medskip

Sin embargo, presenta una desventaja importante: cuando el gradiente es igual a cero, 
los pesos y sesgos no se actualizan durante la retropropagación. Este fenómeno se conoce 
como \guillemotleft ReLU moribundo\guillemotright.
\medskip

Este problema puede mitigarse mediante la función de activación Leaky ReLU, la cual 
permite mantener un gradiente distinto de cero para valores negativos al aplicar 
un factor pequeño $\alpha$, típicamente cercano a $0.01$, a dichos valores. La función Leaky
ReLU está dada por:
\begin{equation}
    f(z) = \begin{cases}
        z & \text{si } z \ge 0\\
        \alpha z & \text{si } z < 0
    \end{cases}
\end{equation}

\subsubsection{Función Softmax}
La función softmax es una generalización de la función sigmoide para problemas de
clasificación multiclase, y se utiliza típicamente en la capa de salida de redes neuronales.
\medskip

La función softmax toma como entrada un vector de valores reales
$z$
\begin{equation}
    z=\langle z_{1}, z_{2}, \dots, z_{n}\rangle
\end{equation} y los transforma en un vector de
probabilidades $p$
\begin{equation}
    p=\langle p_{1}, p_{2}, \dots, p_{n}\rangle
\end{equation}
para $n$ clases, definido como:
\begin{equation}
    p_{i} = \frac{e^{z_{i}}}{\sum_{j=1}^{n} e^{z_{j}}}
\end{equation}
donde $z_i$ corresponde a la entrada neta (pre-activación) de la $i$-ésima neurona en la capa de salida.
\medskip

De esta forma, cada $p_i$ representa la probabilidad de pertenencia a la clase $i$, y la suma
de todas las probabilidades es igual a 1.
\medskip

A diferencia de otras funciones de activación, la función softmax no opera de manera
independiente sobre cada neurona, sino que se aplica al vector completo de salidas de
la capa, ya que cada valor de salida depende de todas las entradas.
% Mencionar linearidad y no-linearidad
% Ejemplos: Sigmoide, ReLU, tanh
\medskip

En conjunto, las funciones de activación desempeñan un papel fundamental en las redes neuronales,
ya que introducen no linealidad en el modelo, permitiendo la representación y el aprendizaje de
patrones complejos. A lo largo de su evolución, se han propuesto distintas funciones con el objetivo
de mejorar la capacidad de entrenamiento y mitigar problemas como el desvanecimiento del gradiente.
En la práctica, funciones como ReLU y sus variantes se emplean comúnmente en las capas ocultas,
mientras que funciones como la sigmoide o softmax se utilizan en la capa de salida dependiendo del
tipo de problema. La elección adecuada de estas funciones influye directamente en el comportamiento
y desempeño de la red, lo cual motiva el estudio de la estructura y organización de las neuronas
dentro de una arquitectura de red neuronal.

\subsection{Arquitectura de Redes Neuronales}
Las redes neuronales artificiales se organizan típicamente en una estructura por capas,
donde cada capa está compuesta por un conjunto de neuronas que procesan la información
de manera conjunta. En este esquema, la salida de una capa sirve como entrada para la
siguiente, formando una composición de transformaciones sucesivas.
\medskip

Generalmente, una red neuronal se divide en tres tipos de capas:
\begin{description}
    \item[Capa de entrada] Esta capa es responsable de recibir los datos o señales de entrada,
    usualmente preprocesarlos (por ejemplo, mediante normalización), y distribuirlos a las
    capas subsecuentes.
    
    \item[Capas ocultas] Estas capas están compuestas por neuronas encargadas de extraer
    características y patrones, y realizan la mayor parte del procesamiento interno de la red.
    
    \item[Capa de salida] Esta capa también está compuesta por neuronas y es responsable
    de producir la salida final de la red con base en el procesamiento realizado por las
    capas anteriores.
\end{description}
\medskip

Las redes neuronales pueden clasificarse en distintas arquitecturas de acuerdo con la
disposición de las neuronas y sus interconexiones.
\medskip

Las redes \textit{feedforward} son aquellas en las que la información fluye en una sola
dirección, desde la capa de entrada hasta la capa de salida, sin ciclos ni retroalimentación.
En este tipo de redes, las neuronas de una capa se conectan típicamente con las neuronas
de la siguiente. Sin embargo, algunas arquitecturas permiten \textit{conexiones de atajo},
en las que las conexiones pueden saltarse una o más capas. Adicionalmente, es común que
las neuronas de una capa se conecten con todas las neuronas de la siguiente, en cuyo caso
se denomina \textit{capa totalmente conectada}.
\medskip

Una red neuronal feedforward puede estar compuesta por una sola capa (sin capas ocultas),
lo que se conoce como arquitectura de una sola capa. Un ejemplo de esta arquitectura es
el perceptrón, uno de los primeros modelos de redes neuronales.
\medskip

Por otro lado, las redes recurrentes son aquellas en las que existen conexiones cíclicas,
permitiendo que la información se retroalimente dentro de la red. Esto posibilita modelar
dependencias temporales o secuenciales. La recurrencia puede presentarse de forma directa,
cuando una neurona se conecta consigo misma, o de forma indirecta, mediante conexiones
a través de otras neuronas.

\begin{figure}[ht]
    \centering
    \includegraphics[width=0.8\textwidth]{diagrams/chapter2/multilayer_nn_arch.pdf} 
    \caption{Red neuronal feedforward con dos capas ocultas totalmente conectadas.}
    \label{fig:rn_arquitectura}
\end{figure}
% Capas: entrada, oculta, salida
% Feedforward flow
% Diagrama

\subsection{Proceso de Aprendizaje}
Una de las características más importantes de las redes neuronales es su capacidad
de generalizar, es decir, de resolver problemas a partir del aprendizaje de patrones
que describen el comportamiento de los datos. Estos patrones se obtienen a partir
de un \textit{conjunto de entrenamiento}.
\medskip

Existen tres paradigmas de aprendizaje principales que definen el proceso general
de aprendizaje: supervisado, no supervisado y por refuerzo.

\begin{description}
    \item[Aprendizaje supervisado] Consiste en el uso de un conjunto de entrenamiento
    en el cual cada patrón de entrada está asociado a una salida correcta. El entrenamiento
    ajusta los pesos de la red considerando la discrepancia entre la salida producida
    por la red y el valor deseado. Se considera que la red ha sido entrenada cuando
    dicha discrepancia se encuentra dentro de un rango aceptable.

    \item[Aprendizaje no supervisado] Este paradigma no requiere conocimiento de las
    salidas deseadas. En su lugar, la red se ajusta en función de similitudes y estructuras
    presentes en los datos de entrada, descubriendo patrones de manera autónoma para
    tareas como agrupamiento o reducción de dimensionalidad.

    \item[Aprendizaje por refuerzo] En este paradigma, un agente interactúa con un entorno
    y recibe señales de recompensa en función de sus acciones. A diferencia del aprendizaje
    supervisado, no se dispone de una salida correcta explícita. En su lugar, el modelo
    aprende a tomar decisiones que maximizan la recompensa acumulada, reforzando
    comportamientos deseables y penalizando aquellos que no lo son.
\end{description}

En adición a los paradigmas de aprendizaje, el proceso de entrenamiento también
puede clasificarse en función de cómo se procesa el conjunto de entrenamiento.
En \textit{offline learning}, también conocido como aprendizaje por lotes, el
ajuste de los pesos se realiza después de procesar un conjunto de muestras de
entrenamiento.
\medskip

Por otro lado, en \textit{online learning}, el ajuste de los pesos se realiza
después de procesar cada muestra de entrenamiento.

% Aprendizaje supervisado vs no supervisado
% Descenso de gradiente
% Funcion de perdida
% Backpropagation
\medskip


El proceso de aprendizaje consiste en la aplicación de un algoritmo que ajusta los
pesos de las conexiones neuronales y los sesgos de las neuronas. El proceso de entrenamiento en el aprendizaje supervisado puede describirse en dos
etapas principales. En la primera etapa, conocida como \textit{forward propagation},
las señales correspondientes a una muestra de entrenamiento se introducen en la red
y se propagan capa a capa hasta producir una salida.
\medskip

A partir de esta salida, se calcula un error en la \textit{computación del error} al compararla
 con el resultado deseado
mediante una función de pérdida, la cual cuantifica el error para una muestra
individual. Una de las funciones de pérdida más utilizadas es la función del error cuadrado,
definida como:

\begin{equation}
    E(k)=\frac{1}{2}\sum_{j=1}^{n}(t_{j}(k)-a_{j}(k))^{2}
\end{equation}

donde:

\begin{description}
    \item[$k$] es la $k$-ésima muestra de entrenamiento.
    \item[$n$] es el número de neuronas en la capa de salida.
    \item[$t_{j}(k)$] es el valor esperado de la $j$-ésima neurona para la $k$-ésima
    muestra de entrenamiento.
    \item[$a_{j}(k)$] es el valor obtenido de la $j$-ésima neurona para la $k$-ésima
    muestra de entrenamiento.
\end{description}

A partir de estas pérdidas se define una función de costo, que agrega
el error sobre el conjunto de entrenamiento, siendo su minimización el principal
objetivo del proceso de aprendizaje. Una de las funciones más utilizadas es el
promedio del error cuadrado, definido como:

\begin{equation}
    E = \frac{1}{m}\sum_{k=1}^{m}E(k)
\end{equation}

donde $m$ es el número de muestras de entrenamiento en el conjunto.
\medskip

El proceso de aprendizaje busca minimizar este costo. 
Esto puede lograrse mediante el descenso por gradiente, un proceso de optimización que busca 
minimizar una función definida en un espacio $n$-dimensional.
El gradiente $\nabla f$ de una función escalar $f$ es un vector que, en cada punto
diferenciable, apunta en la dirección de mayor incremento de la función.
\medskip

De manera análoga, el gradiente negativo $-\nabla f$ indica la dirección de mayor
descenso. Por lo tanto, el descenso por gradiente consiste en actualizar los parámetros
en la dirección del gradiente negativo con el fin de minimizar la función de error.
\medskip

La retropropagación es la siguiente etapa y es un proceso que permite calcular de manera eficiente el
gradiente de la función de costo $\nabla E(W)$ con respecto a los parámetros de la red,
aprovechándose de la regla de la cadena para el cálculo de derivadas. Dado que
la salida de la red puede interpretarse como una composición de funciones a través
de sus capas, es posible encadenar las derivadas desde la capa de salida hacia las
capas anteriores, propagando el gradiente a lo largo de la red.
\medskip

Una vez calculado el gradiente de $E(W)$, se determina el cambio a aplicar a los pesos,
definido como:
\begin{equation}
    \Delta W = -\eta \nabla E(W)
\end{equation}
donde $\eta$ es la tasa de aprendizaje de la red y $\Delta W$ representa el vector
de actualizaciones para todos los pesos.
\medskip

Finalmente, la etapa de optimización actualiza los parametros, sumando el
cambio obtenido a los valores actuales de los parámetros.

\subsection{Limitaciones de Redes Neuronales}
A pesar de la versatilidad de las redes neuronales para modelar y generalizar problemas,
estas presentan ciertas limitaciones que pueden dificultar el proceso de aprendizaje y
predicción en determinados contextos.
\medskip

Una de las principales limitaciones surge al utilizar capas totalmente conectadas.
Dependiendo del número de neuronas de entrada, así como del número de neuronas y capas
ocultas, el modelo puede requerir un gran número de parámetros. Esto incrementa el costo
computacional tanto del entrenamiento como de la inferencia.

Este problema se acentúa especialmente al trabajar con datos de alta dimensionalidad,
como imágenes, donde cada píxel puede representar una entrada. En estos casos, una red
totalmente conectada requeriría un número muy elevado de conexiones entre capas,
lo que hace necesario escalar considerablemente el tamaño del modelo.
\medskip

El aumento en el número de parámetros también incrementa el riesgo de sobreajuste
(\textit{overfitting}), fenómeno que ocurre cuando el modelo se ajusta excesivamente
a los datos de entrenamiento. Como consecuencia, la función de error sobre dichos datos
puede ser muy baja, pero el modelo pierde capacidad de generalización, presentando un
alto error al enfrentarse a datos no vistos previamente.

Esto dificulta el entrenamiento efectivo del modelo y reduce la fiabilidad de sus
predicciones en escenarios reales.
% Muchos parametros
% Falta de espacialidad
% Escalabilidad pobre para imagenes